\begin{table*}[t]
\caption{Main terms of this paper in plain words.}
\label{tab:terms}
\centering\footnotesize
\setlength{\tabcolsep}{4pt}
\renewcommand{\arraystretch}{1.15}
\begin{tabular}{@{}>{\raggedright\arraybackslash}p{32mm}>{\raggedright\arraybackslash}p{123mm}@{}}
\toprule
Term & Meaning in this paper\\
\midrule
Detector & A neural network that finds objects in one video frame and returns a box for each of them.\\
Candidate (detection) & One box returned by the detector, with a class label (for example, car) and a confidence score.\\
Confidence score & A number between 0 and 1 that says how sure the detector is that the box holds an object. Different detectors use different score scales.\\
Threshold & A fixed score limit. The tracker uses a candidate only if its score is at or above the limit.\\
Tracker & A program that links the candidates of consecutive frames into tracks, one track per object.\\
Track, identity (ID) & The boxes of one object over time, with one ID number that should never change.\\
Sequence & One video of a dataset.\\
Stream & The detector output of one video, frame after frame.\\
Host & The tracker that \scl{} serves. It runs after \scl{} with its own published code and settings, which are never changed.\\
Host contract & Four published settings of the host that \scl{} reads (Sect.~\ref{sec:contract}).\\
Matching and birth thresholds & The score a candidate needs to continue an existing track (matching, also called association), and the score it needs to start a new track (birth).\\
Single-stage, two-stage tracker & A single-stage tracker uses one threshold and ignores every candidate below it. A two-stage tracker also uses low-score candidates in a second step, only to continue existing tracks.\\
Frozen & Kept fixed: not retrained, not retuned, not changed in any way.\\
Labels (ground truth) & Boxes and identities drawn by people. They are used only to measure the results, never by \scl{}.\\
Development data & Data that was looked at while \scl{} was designed.\\
Held-out, predeclared & Data or trackers that were chosen in advance, were not used during the design, and were run once with the frozen \scl{}.\\
Logit & The score written on another scale, $\log\bigl(s/(1-s)\bigr)$, Eq.~\eqref{eq:logit}. It spreads out scores close to 0 and 1.\\
Otsu's method & A standard image-processing rule that picks the threshold that best splits a histogram into two groups \cite{otsu1979}.\\
HOTA, MOTA, IDF1 & Tracking metrics in percent; higher is better (Table~\ref{tab:metrics}). HOTA, the main metric, balances finding the objects and keeping their identities.\\
FP, FN, IDS & False positives (boxes where there is no object), false negatives (missed objects), and identity switches (an object's ID changes). They are counts; lower is better.\\
Paired comparison & The same tracker runs twice on the same sequences and detections, once alone and once with \scl{}.\\
95\% interval, significant & The range of likely values of a difference, found by resampling the test sequences 10,000 times (bootstrap). A difference is called significant when this range does not include zero.\\
\bottomrule
\end{tabular}
\end{table*}
